% L3-3 (2026-10-30): new subsubsection for §6.5 of numerical_results.tex (label subsubsec:num_L3_3), figure environment included.
% Figure: plots/L3-3/single_resonator.png -> figures/L3_3_single_resonator.png (scripts/export_paper_assets.jl L3-3).
% Data: data/exp_pro/L3-3/{frequencies,zeta_sweep,sensitivity,renormalised,checks,fine_mesh_reference}.csv (203/203 checks).
% Matrix row L3.3 (frozen 2026-10-29); post-freeze deviation: h = lambda_p/4.5 instead of lambda_p/5 (20 GiB memory cap).
{\color{red}
\subsubsection{Single resonator on a large plate}\label{subsubsec:num_L3_3}

\paragraph{Test description.} One point resonator ($M_\rho=0.01~\si{\meter\cubed}$, $\omega_r=10~\si{\radian\per\second}$) is attached at the centre of a plate that covers the whole free surface of a fluid of depth $H=10~\si{\meter}$, with $C_\rho=0.49~\si{\meter\tothe{5}\per\second\squared}$, $m_\rho=0.01~\si{\meter}$ and $\nu=0$. For this host, $\gamma_r=k_0/(2\mathcal{D}_0')=0.0255$ at $\omega=10~\si{\radian\per\second}$, the renormalised resonance of the theory is $\tilde\omega=9.9506~\si{\radian\per\second}$ and critical coupling is reached at $\zeta_\mathrm{crit}=0.0129$. Because the resonator and its point load are axisymmetric, we solve one quarter of a square of side $16\lambda_p$. On the symmetry planes the fluid has no flux and the plate has $\partial_n\eta=0$, imposed by Nitsche terms that are the restriction of the interior C/DG facet terms of the mirrored plate. An absorbing frame of Kelvin--Voigt damping starts at $4.9\lambda_p$ from the centre and is $3.1\lambda_p$ wide. The fields are $Q_3$, with horizontal cells of $\lambda_p/4.5$ and eight graded vertical cells ($N=3.1\times10^5$). We solve the bare plate under a unit point load at the attachment point, one factorisation per frequency. This gives $G_h(\bm 0)$ and the radiated field $G_h(\xvec)$. By linearity, the force on the resonator under an incident wave of amplitude $A$ is $F_1/A=1/(\alpha^{-1}-G_h(\bm 0))$ and the scattered field is $F_1G_h(\xvec)$, for every damping ratio $\zeta$ (Sherman--Morrison on the same factorisation). We fit the outgoing amplitude $C^+$ of $G_h$ on rings $1.5\lambda_p\le r\le4.4\lambda_p$, at 19 angles over the quarter, with $H_0^{(1)}(kr)$, $H_0^{(2)}(kr)$ and the first evanescent modes. From $C^+$ we obtain $\sigma_\mathrm{sc}=4|F_1C^+/A|^2/k_0$ and the extinction $\sigma_\mathrm{ext}=-(4/k_0)\,\mathrm{Re}(F_1C^+/A)$, and $\sigma_\mathrm{abs}=2\gamma_d|F_1/A|^2/\mathcal{D}_0'$ from the dissipated power. Eleven frequencies cover $\tilde\omega\pm0.56~\si{\radian\per\second}$, i.e.\ three radiative half-widths on each side. We also run a shorter domain, with the frame starting at $4\lambda_p$, and a coarser mesh, $\lambda_p/4$. Identities are checked on a small mesh against the full domain and against the resonator solved inside the finite element system. Every run stays below $16$~GiB of memory.

\paragraph{Results.} The quarter-domain solution agrees with the full domain to $4.5\times10^{-13}$, and the Sherman--Morrison force agrees with the monolithic resonator to $7\times10^{-14}$. At all eleven frequencies the radiation of the point load matches the theory. $\mathrm{Im}\,G_h(\bm 0)$ equals $\gamma_r$ within $0.38\%$ and the fitted amplitude gives $|C^+/(\mathrm{i}\gamma_r)-1|\le1.0\times10^{-3}$. The radiated power through the rings equals the work of the point load within $0.29\%$, the reflection from the frame is below $9.3\times10^{-4}$, and the fitted wavenumber is within $0.04\%$ of $k_0$. The far field is isotropic within $1.7\%$ over the quarter. $\mathrm{Re}\,G_h(\bm 0)$ is $2.8$--$3.9\%$ below the contour-integral value, since the driving point depends on the near field and the mesh is coarse there. The renormalised resonance of the finite element model is $\tilde\omega_h=9.9522~\si{\radian\per\second}$, within $1.6\times10^{-4}$ of $\tilde\omega$, and its optimal damping $\zeta_{\mathrm{opt},h}=0.01293$ is within $0.3\%$ of $\zeta_\mathrm{crit}$. The absorption cross-section never exceeds $1/k_0$: over all frequencies and $83$ damping ratios between $10^{-3}$ and $10^{-1}$, the maximum of $k_0\sigma_\mathrm{abs}$ is $0.996$ (Figure~\ref{fig:num_L3_3}). At critical coupling the bound is attained, $k_0\sigma_\mathrm{abs}=\gamma_r/\mathrm{Im}\,G_h=0.997$ at $(\tilde\omega_h,\zeta_{\mathrm{opt},h})$, and absorbed and scattered powers are equal within $0.26\%$. The optical theorem $\sigma_\mathrm{ext}=\sigma_\mathrm{sc}+\sigma_\mathrm{abs}$ holds within $0.32\%$ at $\zeta=0$, $\zeta_\mathrm{crit}$ and $3\zeta_\mathrm{crit}$ for every frequency. $k_0\sigma_\mathrm{abs}$ agrees with the single-resonator theory within $0.006$ and $\sigma_\mathrm{sc}$ within $1.4\%$. Moving the frame inwards by $0.9\lambda_p$ changes $G_h(\bm 0)$ by $0.20\%$, and the coarser mesh changes it by $0.46\%$. We conclude that the finite element model of a single resonator on a three-dimensional plate reproduces the absorption bound $\sigma_\mathrm{abs}\le1/k_0$, its attainment at critical coupling and the optical theorem.

\begin{figure}[t]
    \centering
    \includegraphics[width=\linewidth]{figures/L3_3_single_resonator.png}
    \caption{Single resonator at the centre of a large plate (quarter domain, $h=\lambda_p/4.5$). (a) $k_0\sigma_\mathrm{abs}$ against $\omega$ at $\zeta_\mathrm{crit}$ and $3\zeta_\mathrm{crit}$, and $k_0\sigma_\mathrm{sc}/4$ at $\zeta=0$; lines: single-resonator theory; markers: finite element; dotted: the bound $1/k_0$. (b) $k_0\sigma_\mathrm{abs}$ and $k_0(\sigma_\mathrm{ext}-\sigma_\mathrm{sc})$ against $\zeta$ at $\tilde\omega$, all from one factorisation. (c) $|G_h(x,0)|/\gamma_r$ along the axis and the diagonal at $\tilde\omega$, and the infinite-plate Green function; grey: absorbing frame.}
    \label{fig:num_L3_3}
\end{figure}
}
